\chapter{Theoretische Grundlagen}

Das übergreifende Thema war die objektorientierte Programmierung. Dies bedeutet, dass man verschiedene Klassen anlegt, die anschließend durch programmierte Methoden und Funktionen miteinander interagieren können.\\
Dadurch kann man einerseits größere Projekte realisieren, indem man parallel an verschiedenen Klassen arbeiten kann, andererseits wird dadurch mehr Übersichtlichkeit geschaffen, weil man verschiedene, thematisch geordnete, Dateien erstellt.\\
Eine Klasse ist eine Datei, in der ein Objekt oder eine Szene mit ihren Parametern und Methoden und Funktionen definiert wird\\
Eine Methode ist Code, der einen oder mehrere Parameter bekommt und dann gewisse, vorher festgelegte Operationen durchführt.\\
Eine Funktion ist Code, der einen oder mehrere Parameter bekommt und daraufhin Operationen durchführt und abschließend einen Wert zurück gibt.\\
Ein Objekt ist eine Klasse, die Methoden, Funktionen und einen Konstruktor enthält.\\
Eine Szene ist eine Klasse, die Methoden, Funktionen und ein Bild enthält, welches angezeigt wird, wenn die Szene aufgerufen wird.\\
Ein Konstruktor ist Code, in dem eine Instanz eines Objekts "konstruiert" wird.\\
Eine Instanz ist ein "Körper" der nach den Regeln seiner dazugehörigen Klasse funktioniert. Eine Klasse ist somit nur das Regelwerk während die Instanz die ausführende Gewalt ist.\\

Wir haben uns für Java entschieden, weil es die Programmiersprache ist, mit der die meisten der Gruppe angefangen haben zu programmieren und somit auch am besten beherrschen. Außerdem ist Java auch die Programmiersprache, die meistens für objektorientierte Programmierung benutzt wird.\\